\chapter{Type 1 Diabetes Mellitus}
\index{Type 1 Diabetes Mellitus}
\label{sec:Type 1 Diabetes Mellitus}

\section{Overview}
Type 1 diabetes mellitus (T1DM) is an autoimmune disease characterized by destruction of pancreatic beta cells in the islets of Langerhans, leading to absolute insulin deficiency. It accounts for 5--10\% of diabetes cases. This metabolic disorder involves dysregulation of normal bodily processes, often requiring ongoing monitoring and management. It is increasingly common due to lifestyle and dietary factors. Metabolic disorders involve disruption of normal biochemical processes essential for health. These conditions may result from enzyme deficiencies, hormonal imbalances, or lifestyle factors. Many metabolic disorders are chronic and require ongoing management. Lifestyle modifications are often the cornerstone of treatment.
\section{Causes}
The underlying mechanisms involve complex interactions between genetic predisposition and environmental factors. Disruption of normal physiological processes leads to the characteristic features of the condition. Multiple contributing factors may act synergistically to trigger or worsen the condition.

\begin{itemize}[noitemsep]
  \item This condition happens because of autoreactive T cells (with CD8+ cytotoxic T cells being primary effectors) that destroy beta cells, with autoantibodies (anti-GAD65, anti-IA-2, anti-insulin, anti-ZnT8) serving as biomarkers. Genetic susceptibility is linked to HLA-DR3/DR4 and HLA-DQ2/DQ8 haplotypes
  \item Environmental triggers (viral infections, early cow's milk exposure) may initiate autoimmunity in predisposed individuals.
\end{itemize}
\section{Risk Factors}
\begin{itemize}[noitemsep]
  \item Genetic predisposition and family history
  \item Advancing age
  \item Lifestyle factors (diet, exercise, smoking, alcohol)
  \item Environmental exposures
  \item Underlying medical conditions
  \item Medication use
\end{itemize}
\section{Signs and Symptoms}

\subsection{Common symptoms}
\begin{itemize}[noitemsep]
  \item You may experience the classic triad of polyuria, polydipsia, and weight loss, with diabetic ketoacidosis (DKA) as the initial presentation in 30--40\% of children.
  \item Pain or discomfort in the affected area
  \end{itemize}

\subsection{Emergency warning signs}
\begin{itemize}[noitemsep]
  \item Severe or worsening symptoms of type 1 diabetes mellitus require immediate medical evaluation
\end{itemize}
\section{Progression and Stages}
The condition may progress through different stages of severity over time. Early stages often involve mild or subtle symptoms that may be overlooked. Moderate stages involve more noticeable symptoms affecting daily function. Advanced stages may involve significant impairment and complications.
\section{Complications}
\begin{itemize}[noitemsep]
  \item Your outlook depends on glycemic control
  \item Microvascular and macrovascular complications determine long-term outcomes.
  \item Disease progression leading to worsening symptoms
  \item Secondary infections or injuries
  \item Side effects of treatment
  \item Reduced quality of life
  \item Emotional and psychological impact
\end{itemize}
\section{Diagnosis}
To make the diagnosis, doctors look for fasting plasma glucose $\geq$126 mg/dL, random glucose $\geq$200 mg/dL with symptoms, HbA1c $\geq$6.5\%, or 2-hour OGTT glucose $\geq$200 mg/dL.

\begin{itemize}[noitemsep]
  \item Comprehensive medical history and physical examination
  \item Laboratory tests to assess organ function
  \item Imaging studies to visualize affected structures
  \item Specialized testing depending on the affected system
  \item Consultation with relevant specialists
\end{itemize}
\section{Prevention}
\begin{itemize}[noitemsep]
  \item Maintain a healthy lifestyle with balanced diet and exercise
  \item Avoid known risk factors
  \item Attend regular medical check-ups for early detection
  \item Follow recommended screening guidelines
\end{itemize}
\section{Treatment Options}
\begin{itemize}[noitemsep]
  \item Treatment typically involves lifelong insulin therapy via multiple daily injections (basal-bolus regimen) or continuous subcutaneous insulin pump
  \item Continuous glucose monitoring (CGM) and insulin pump systems (closed-loop "artificial pancreas" systems) improve glycemic control.
  \item Medications to manage symptoms and address underlying mechanisms
  \item Lifestyle modifications and self-care measures
  \item Physical therapy and rehabilitation
  \item Surgical intervention when indicated
  \item Regular monitoring and follow-up care
\end{itemize}
\section{Living with It}
\begin{itemize}[noitemsep]
  \item Follow your treatment plan as prescribed
  \item Maintain regular follow-up with your healthcare provider
  \item Make necessary lifestyle adjustments
  \item Seek support from family, friends, or support groups
  \item Stay informed about your condition
\end{itemize}
\section{Outlook}
Islet transplantation and Teplizumab (anti-CD3 antibody) are emerging therapies.
